\documentclass[10pt,a4paper]{amsart}
\usepackage[margin=19mm]{geometry}
\usepackage{amsmath,amssymb,mathtools}
\usepackage[scr=rsfs,cal=euler]{mathalfa}
\usepackage{xfrac}
\usepackage[unicode,hidelinks,bookmarks=false]{hyperref}
\newcommand{\ie}{i.e.\ }
\newcommand{\cf}{cf.\ }
\newcommand{\resp}{respectively\ }
\newcommand{\Subsection}{\subsection}
\providecommand{\nobreakdash}{\nobreak\hbox{-}}
\setlength{\emergencystretch}{2em}
\multlinegap0pt
\usepackage{bm}
\usepackage{bbm}
\usepackage{dsfont}
\usepackage{xspace}
\usepackage{cancel}
\usepackage{mathtools}
\usepackage{amsthm}
\makeatletter
\@ifundefined{theorem}{\newtheorem{theorem}{Theorem}}{}
\@ifundefined{lemma}{\newtheorem{lemma}[theorem]{Lemma}}{}
\makeatother
\title[The weight two and opposite sign cases for the Fourier relat]{The weight two and opposite sign cases for the Fourier relative trace formulas}
\author{Matteo Di Scipio}
\date{Source: arXiv:2603.13062v2 (CC BY 4.0)}
\begin{document}
\maketitle

\noindent\emph{Source and licence.} The weight two and opposite sign cases for the Fourier relative trace formulas. Version arXiv:2603.13062v2 (CC BY 4.0), distributed under the Creative Commons Attribution 4.0 International licence, as recorded in the arXiv licence metadata for this exact version and shown on the abstract page https://arxiv.org/abs/2603.13062v2. Full rights notice: licenses/arxiv-2603.13062-CC-BY-4.0.txt.\par\smallskip

\noindent\textit{Editorial scope.} Counted unit: the Lemma whose source label is \texttt{lemma:derivativebound} in the article (source0.tex lines 921--926, source label \texttt{lemma:derivativebound}), delivered together with the article's own complete original proof of it (source0.tex lines 927--951, one \texttt{proof} environment of 25 lines). No mathematics is written, repaired or re-derived: statement and proof are the article's own text line for line. Required context retained uncounted: none. Every definition, display and label of those lines is retained unchanged. Omitted from this package and not redistributed: 1--920, 952--1145. Nothing inside a retained excerpt is omitted or altered: every retained body is delivered exactly as the article's lines are written, so no omission receipt of any kind accompanies this package. Citations: the retained excerpts carry no citation command at all, so no bibliography entry of the article is transcribed and no reference is redistributed. Reference adaptation: none was needed and none was made; every reference inside the delivered unit resolves inside it, so no unresolved reference or citation is produced. Printed slips kept literal: no printed word, symbol, hypothesis or number of the counted statement or of its proof was altered. Layout and class: the article's own class is \texttt{article} with packages including biblatex, graphicx, mathrsfs, geometry, mathtools, enumitem, amsmath, amsthm, amssymb, xcolor, hyperref, sectsty; the wrapper is the lane's pinned \texttt{amsart} editorial wrapper at 10pt on A4, which restates the theorem-like environments the retained text needs and adds the article's own macro definitions that the retained text needs (none), with extra wrapper packages (bm, bbm, dsfont, xspace, cancel, mathtools). The delivered numbering is the unit's own. Assets: no retained span contains a figure environment, a graphics-inclusion command, a PDF-page inclusion, a table or a TikZ picture; no image file is copied into this package, the package declares no asset dependency and no image file is redistributed with it. Licence: version arXiv:2603.13062v2 of this article is distributed under the Creative Commons Attribution 4.0 International licence, as recorded in the arXiv licence metadata for this exact version and shown on the abstract page https://arxiv.org/abs/2603.13062v2. A full rights notice is delivered at licenses/arxiv-2603.13062-CC-BY-4.0.txt. Characters: every retained excerpt is pure ASCII, so the delivered engine, XeLaTeX with its default Latin Modern text font, reports no missing character.
\begin{lemma}\label{lemma:derivativebound}
    Let $m\geq 1$ be a positive integer, then for $T\geq 10$ say,
    \begin{equation}\Bigr|\frac{\partial^m}{\partial y^m}\rho(r_T(z))\Bigr|\ll y^{-m}
    \end{equation}
    where $z=x+iy \in\mathbb{H} $ and the implicit constant depends only on $m$ and $\rho$.
\end{lemma}

\begin{proof}
    For sake of notation, we take $r^{(m)}(z)$ to denote the $m$-th derivative with respect to $y$. By the generalised chain rule 
    $$\frac{\partial^m}{\partial y^m}\rho(r_T(z))=\sum_{j=1}^{m}\sum_{\substack{a_1+...+a_j=m \\m\geq a_1\geq ... \geq a_j \geq 1 }}A_{a_1,...,a_j}\rho^{(m)}(r(z))\ r^{(a_1)}(z)...r^{(a_j)}(z)$$
    where $A_{a_1,...a_j}$ are some non-negative constants. We claim that for $a>0$
    $$r^{(a)}(z)=\frac{P_a(x,y)}{y^a((x^2+y^2+1)^2-4y^2)^{a-1/2}}$$
    with $P_a(z)$ a polynomial of total degree $\leq 4a-2$.
    As the base case
    $$r^{(1)}(z)=\frac{y^2-x^2-1}{y((x^2+y^2+1)^2-4y^2)^{1/2}}$$
    So suppose our claim holds for some $a$, then by the product rule
    \begin{align*}
            r^{(a+1)}(z)  = & -\frac{aP_a(x,y)}{y^{a+1}((x^2+y^2+1)^2-4y^2)^{a-1/2}}+\frac{P^{(1)}(x,y)}{y^{a}((x^2+y^2+1)^2-4y^2)^{a-1/2}} \\
            & \hspace{1cm} -\frac{2(2a-1)y(y^2+x^2-1)P_a(x,y)}{y^a((x^2+y^2+1)^2-4y^2)^{a+1/2}}
    \end{align*}
    Collecting terms, the denominator has the desired form, and the numerator is
    \begin{align*}
        -a((x^2+y^2+1)^2-4y^2)P_{a}(x,y)+y((x^2+y^2+1)^2-4y^2)& P_a^{(1)}(x,y) \\ 
        & -2(2a-1)y^2(x^2+y^2-1)P_a(x,y)
    \end{align*}
    whose degree satisfies the desired bound, thus the claim.\\
    Rearranging
    $$r^{(a)}(z)=\frac{1}{y^a}\frac{P_a(x,y)}{(x^2+y^2+1)^{2a-1}}\left(1-\frac{4y^2}{(x^2+y^2+1)^2}\right)^{1/2-a}$$
    By the support of the derivatives of $\rho$, we may assume $x+iy$ lies outside a small neighborhood of $i$, thus the third factor is $\ll_a 1$. For the second factor, if $|x|\geq y$ then $|P_a(x,y)|\ll_a |x|^{4a-2}$ and 
    $$\frac{|x|^{4a-2}}{(x^2+1)^{2a-1}}\leq 1$$
    A similar result holds for $|x| \leq y$, thus proving the lemma.
\end{proof}

\end{document}
